
D\'
emontrons que (i) implique (ii), (iii), et (iv). Supposons donc que \(\mathcal L\) soit limit\'ee par un faisceau inversible sur un \(X_T\). Alors \(\mathcal L\) satisfait \`a (iv) par EGA II 4.6.12. Supposons \(L\) plat sur \(T\). Alors le polyn\^ome en \(m\) et \(n\), \(\chi(L_t^{\otimes m}(n))\) ne d\'epend que de la composante connexe de \(T\) (EGA III 7.9.4). Par suite \(\mathcal L\) satisfait \`a (ii), et compte tenu de 2.1, \`a (iii).

Notons ensuite que la famille des classes des \(H_K\) (resp. \({O}_{X_K}\)) est limit\'ee, donc satisfait \`a (iii) (resp. \`a (ii)).

Supposons que \(\mathcal L\) satisfasse \`a (iv). Alors \(\langle [N\cdot c_1(H_K)+c_1(L_K)]\cdot c_1(H_K)^{r-1}\rangle>0\), d'o\`u \(d(L_K)>-Nd(H_K)\), et de m\^eme, \(Nd(H_K)>d(L_K)\). De plus, \(\langle [N c_1(H_K)+c_1(L_K)]^2\cdot c_1(H_K)^{r-2}\rangle>0\), d'o\`u \(\langle c_1(L_K)^2\cdot c_1(H_K)^{r-2}\rangle>-2Nd(L_K)-N^2d(H_K)\). Donc \(\mathcal L\) satisfait \`a (iii).

D\'emontrons enfin que (ii) (resp. (iii)) implique (i). Gr\^ace \`a 2.10 et 2.3, on peut remplacer \(H\) par un multiple\,; supposons donc \(H\) tr\`es ample. \'Evidemment on peut couper \(S\) en morceaux\,; supposons donc qu'il existe une section lin\'eaire \(Y\) de \(X\), qui est int\`egre sur \(S\), \`a fibres de dimension \(2\). En vertu de 3.8, on peut remplacer \(X,\mathcal L\) par \(Y,\mathcal L|Y\), (compte tenu de 1.7 (resp. la formule de projection), qui implique que la condition (ii) (resp. (iii)) sur \(X,\mathcal L\) est \'equivalente \`a (ii) (resp. \`a (iii)) sur \(Y,\mathcal L|Y\))\,; supposons donc \(r=2\). Si \(\delta=\sup\{d(L_K)\}\) est \(>0\), on peut remplacer \(\mathcal L\) par la famille des classes des \(L_K(-\delta)\) gr\^ace \`a 2.3 et 2.10\,; supposons donc \(\delta\le0\).

Si dans les polyn\^omes de Hilbert \(\sum_{i=0}^{2}a_i\binom{n+i}{i}\), \(a_1\) reste born\'e et \(a_0\) reste minor\'e, alors les trois coefficients restent born\'es gr\^ace \`a 2.3 et 2.7\,; donc \(\mathcal L\) est limit\'ee en vertu de 2.11\,; ainsi (ii) implique
